\documentclass[10pt,a4paper]{amsart}
\usepackage[margin=19mm]{geometry}
\usepackage{amsmath,amssymb,mathtools}
\usepackage[scr=rsfs,cal=euler]{mathalfa}
\usepackage{xfrac}
\usepackage[unicode,hidelinks,bookmarks=false]{hyperref}
\newcommand{\ie}{i.e.\ }
\newcommand{\cf}{cf.\ }
\newcommand{\resp}{respectively\ }
\newcommand{\Subsection}{\subsection}
\providecommand{\nobreakdash}{\nobreak\hbox{-}}
\setlength{\emergencystretch}{2em}
\multlinegap0pt
\usepackage{amssymb}
\usepackage{mathtools}
\usepackage{xcolor}
\newtheorem{theorem}{Theorem}
\newtheorem{lemma}{Lemma}
\newcommand{\RR}{\mathbb R}
\newcommand{\Rrate}{\mathcal{R}}
\newcommand{\ev}{\mathrm{ev}}
\newcommand{\norm}[1]{\lVert#1\rVert}
\newcommand{\sph}{\mathbb S}
\newcommand{\vs}{\mathbf{s}}
\newcommand{\vx}{\mathbf{x}}
\newcommand{\vy}{\mathbf{y}}
\title{A Note on Sphere Packing Bounds for Tuple Lattice Sieving}
\author{Thijs Laarhoven}
\date{Source: arXiv:2609.08190v1 (CC BY 4.0)}
\begin{document}
\maketitle

\noindent\emph{Source and licence.} A Note on Sphere Packing Bounds for Tuple Lattice Sieving. Version arXiv:2609.08190v1 (CC BY 4.0), distributed by arXiv under the Creative Commons Attribution 4.0 International licence, http://creativecommons.org/licenses/by/4.0/, as recorded in the arXiv licence metadata for this exact version and shown on the abstract page https://arxiv.org/abs/2609.08190v1. Full licence notice: \texttt{licenses/arxiv-2609.08190-CC-BY-4.0.txt}.\par\smallskip

\noindent\textit{Editorial scope.} Counted unit: the article's theorem thm:upper (source0.tex lines 159--172). The article's complete original proof of it is delivered with it (source0.tex lines 223--235, one \texttt{proof} environment of 13 lines, which the article places away from the statement). No mathematics is written, repaired or re-derived: the statement and the proof are the article's own lines, in the article's own notation, and no retained line is substituted or re-flowed.  Retained uncounted as the source-local context the proof needs: lines 176--182; lines 195--198; lines 208--211.  Nothing was omitted from any retained span, so the package carries no omission record.  No retained line contains a citation, so the wrapper carries no bibliography and no bibliography entry.  No retained line contains a figure environment, an image-inclusion command or drawing code, so no image is delivered and the package declares no asset dependency.  Editorial formatting only, in addition: an AMS \texttt{amsart} preamble that restates the article's own declarations and macros used by the retained lines, namely the environments \texttt{theorem}, \texttt{lemma} on separate counters, together with the standard packages those lines need.  The retained environments print in this document as Theorem 1, Lemma 1 in order of appearance, whereas the article prints the same items with the numbering of its own full document.  The complete native source of the article as distributed in the arXiv e-print for this exact version is bundled unchanged as \texttt{sources/209703/source0.tex}.
\begin{lemma}[A global signing with many short sums]\label{lem:signing}
Let $\vx_1,\ldots,\vx_N\in\sph^{n-1}$ and $1\le r\le N$. There is a signing $\vs=(s_1,\ldots,s_N)\in\{-1,1\}^N$ such that, writing $\vy_A=\sum_{i\in A}s_i\vx_i$,
\begin{align}\label{eq:short-sums}
 \#\left\{A\in\binom{[N]}r:\norm{\vy_A}^2\le r\right\}
 \ge2^{1-r}\binom Nr.
\end{align}
\end{lemma}

\begin{lemma}[Separation of half-tuple sums]
\label{lem:separation}
Suppose $C$ is even-irreducible through order $2m$. For any global signing $\vs$ and any $r\le m$, the sums $\vy_A$ over distinct $r$-subsets obey $\norm{\vy_A-\vy_B}>1$. In particular, all these sums are distinct.
\end{lemma}

\begin{lemma}[Lifting a ball packing to a sphere]
\label{lem:lift}
Let $Y\subset\RR^n$ lie in the closed ball of radius $\rho>0$ centered at the origin, and suppose distinct points of $Y$ have distance greater than $\ell>0$. If $\ell/\rho\le2$, then $|Y|\le A(n+1,1-\ell^2/(2\rho^2))$.
\end{lemma}

\begin{theorem}[Half-tuple packing bound]
\label{thm:upper}
% For $k\ge2$ and $n \geq 2$,
% \begin{align}\label{eq:upper-size-function}
%  N_k(n)\le\min_{1\le r\le\lfloor k/2\rfloor}
%  \left\{r-1+\left[
%   2^{r-1}r!\,A\!\left(n+1,1-\frac1{2r}\right)
%  \right]^{1/r}\right\}.
% \end{align}
For every fixed $k\ge2$,
\begin{align}\label{eq:main-upper}
 \Rrate_k \le \min_{1\le r\le\lfloor k/2\rfloor} \frac{1}{r} \, \kappa\!\left(1-\frac1{2r}\right).
\end{align}
\end{theorem}

\begin{proof}[Proof of Theorem~\ref{thm:upper}]
Put $m=\lfloor k/2\rfloor$. Fix $r\le m$ and a set of size $N\ge r$ that is even-irreducible through order $2m$. Lemmas~\ref{lem:signing} and~\ref{lem:separation} produce at least $2^{1-r}\binom Nr$ distinct, mutually separated points in a ball of radius $\sqrt r$. Lemma~\ref{lem:lift}, with $\rho=\sqrt r$ and $\ell=1$, gives
\begin{align}\label{eq:finite-upper}
 2^{1-r}\binom Nr\le A\!\left(n+1,1-\frac1{2r}\right).
\end{align}
Minimizing over $r$ proves the cardinality bound for even-irreducible sets. Full irreducibility implies these even conditions, so the bound applies to $N_k(n)$ as well.

For each fixed $r$, the same estimate gives
\begin{align}
 \log_2N\le\frac1r\log_2 A\!\left(n+1,1-\frac1{2r}\right)+O_r(1).
\end{align}
Divide by $n$, take the upper limit as $n\to\infty$, and use $(n+1)/n\to1$. The definition of $\kappa$ then gives $\Rrate^{\ev}_{2m}\le r^{-1}\kappa(1-1/(2r))$. Minimizing over $r\le m$ proves the rate statements, including for odd $k$, since $\Rrate_k\le\Rrate^{\ev}_{2m}$.
\end{proof}

\end{document}
